% ============================================================================
% MUSTERLÖSUNG DS 9 - Gruppe A + B: Repaso y Preparación
% ============================================================================
% Spanisch Spätbeginner - Einstimmung vor Beobachtungsstunde
% Lösungen in ROT (#c62828)
% ============================================================================

\documentclass[11pt,a4paper]{article}

\usepackage[utf8]{inputenc}
\usepackage[T1]{fontenc}
\usepackage[ngerman]{babel}
\usepackage[left=2cm,right=2cm,top=1.8cm,bottom=1.5cm]{geometry}
\usepackage{helvet}
\usepackage{xcolor}
\usepackage{tabularx}
\usepackage{array}
\usepackage{enumitem}
\usepackage{tcolorbox}
\usepackage{fancyhdr}
\usepackage{lastpage}

\renewcommand{\familydefault}{\sfdefault}

\definecolor{spanisch}{HTML}{c41e3a}
\definecolor{gruppeB}{HTML}{1565c0}
\definecolor{grau}{HTML}{666666}
\definecolor{hellgrau}{HTML}{f5f5f5}
\definecolor{loesung}{HTML}{c62828}

\pagestyle{fancy}
\fancyhf{}
\fancyhead[L]{\small\textcolor{grau}{Spanisch Spätbeginner -- MUSTERLÖSUNG}}
\fancyhead[R]{\small\textcolor{loesung}{ML-09}}
\fancyfoot[C]{\small\thepage/\pageref{LastPage}}
\renewcommand{\headrulewidth}{0.5pt}

\newcommand{\lsg}[1]{\textcolor{loesung}{\textbf{#1}}}
\newcommand{\es}[1]{\textit{#1}}

\newtcolorbox{gruppenbox}[1][]{
    colback=white,
    colframe=spanisch,
    fonttitle=\bfseries\color{white},
    colbacktitle=spanisch,
    title=#1,
    boxrule=1pt,
    arc=0pt,
    left=5pt, right=5pt, top=5pt, bottom=5pt
}

\newtcolorbox{gruppeBbox}[1][]{
    colback=white,
    colframe=gruppeB,
    fonttitle=\bfseries\color{white},
    colbacktitle=gruppeB,
    title=#1,
    boxrule=1pt,
    arc=0pt,
    left=5pt, right=5pt, top=5pt, bottom=5pt
}

\begin{document}

% ============================================================================
% KOPFBEREICH
% ============================================================================
\begin{center}
{\Large\textbf{\textcolor{loesung}{MUSTERLÖSUNG}}}\\[0.3em]
{\large DS 9 -- Repaso: Describir personas}
\end{center}

\vspace{0.5em}
\hrule
\vspace{1em}

% ============================================================================
% GRUPPE A (A2)
% ============================================================================
\begin{gruppenbox}[Gruppe A (A2): Repaso -- Describir personas]

\textbf{Kasten 1: \es{ser} vs. \es{estar}}

\begin{tabularx}{\textwidth}{@{}|p{2.5cm}|X|X|@{}}
\hline
\textbf{Verb} & \textbf{Verwendung} & \textbf{Beispiel} \\
\hline
\textbf{ser} & \lsg{Eigenschaften} (dauerhaft) & \es{Soy} \lsg{simpático/-a}. \\
\hline
\textbf{estar} & \lsg{Zustände / Gefühle} (vorübergehend) & \es{Estoy} \lsg{cansado/-a}. \\
\hline
\end{tabularx}

\vspace{0.3em}
\textbf{Merke:} \es{ser} = \lsg{wie ich BIN (Charakter)} \hspace{0.5cm} \es{estar} = \lsg{wie ich mich FÜHLE}

\vspace{0.8em}

\textbf{Kasten 2: \es{hay que} + Infinitiv}

Bedeutung: \es{hay que} = \lsg{man muss}

Struktur: \es{hay que} + \lsg{Infinitiv}

Beispiel: \es{Para ser médico hay que ser} \lsg{paciente}.

\vspace{0.8em}

\textbf{Aufgabe 1: \es{ser} oder \es{estar} ergänzen (6 P.)}
\begin{enumerate}[label=\alph*), leftmargin=2em, nosep]
    \item Mi hermano \lsg{es} muy inteligente.
    \item Hoy \lsg{estoy} cansada porque no he dormido bien.
    \item ¿Cómo \lsg{es} tu mejor amiga? -- Es simpática.
    \item Mis padres \lsg{son} muy trabajadores.
    \item ¿Por qué \lsg{estás} triste hoy?
    \item Nosotros \lsg{somos} de Alemania.
\end{enumerate}

\vspace{0.5em}

\textbf{Aufgabe 2: Wortschatz Adjektive (8 P.)}

\begin{tabularx}{\textwidth}{@{}|X|X|X|X|@{}}
\hline
simpático/-a & \lsg{sympathisch, nett} & paciente & \lsg{geduldig} \\
\hline
inteligente & \lsg{intelligent} & divertido/-a & \lsg{lustig, unterhaltsam} \\
\hline
trabajador/-a & \lsg{fleißig} & responsable & \lsg{verantwortungsvoll} \\
\hline
creativo/-a & \lsg{kreativ} & organizado/-a & \lsg{organisiert} \\
\hline
\end{tabularx}

\vspace{0.5em}

\textbf{Aufgabe 3: \es{hay que ser...} (6 P.)}
\begin{enumerate}[label=\alph*), leftmargin=2em, nosep]
    \item Para ser médico/-a, hay que ser \lsg{paciente / responsable}.
    \item Para ser profesor/-a, hay que ser \lsg{paciente / organizado/-a}.
    \item Para ser artista, hay que ser \lsg{creativo/-a}.
\end{enumerate}

\vspace{0.5em}

\textbf{Aufgabe 5: Textproduktion (8 P.) -- Beispiellösung:}

\lsg{Mi mejor amiga se llama María. Tiene 17 años. Es muy simpática y también inteligente. Es trabajadora porque estudia mucho. También es divertida y creativa.}

\vspace{0.3em}
\textit{\textcolor{grau}{Bewertung: 4 Sätze min. (4P), ser/tener korrekt (2P), Konnektoren (2P)}}

\end{gruppenbox}

\vspace{1em}

% ============================================================================
% GRUPPE B (B1)
% ============================================================================
\begin{gruppeBbox}[Gruppe B (B1): Repaso avanzado -- Describir personas]

\textbf{Kasten 1: \es{ser} vs. \es{estar} -- Bedeutungsunterschiede}

\begin{tabularx}{\textwidth}{@{}|p{2.5cm}|X|X|@{}}
\hline
\textbf{Verb} & \textbf{Verwendung} & \textbf{Beispiele} \\
\hline
\textbf{ser} & \lsg{Eigenschaften, Herkunft, Beruf} & \es{Soy inteligente. / Es de Madrid.} \\
\hline
\textbf{estar} & \lsg{Zustände, Gefühle, Orte} & \es{Estoy cansado. / Está en casa.} \\
\hline
\end{tabularx}

\vspace{0.3em}
\textbf{Bedeutungsänderung:}
\begin{tabular}{ll}
\es{ser listo} = \lsg{klug sein} & \es{estar listo} = \lsg{bereit sein} \\
\es{ser aburrido} = \lsg{langweilig sein} & \es{estar aburrido} = \lsg{sich langweilen} \\
\end{tabular}

\vspace{0.5em}

\textbf{Kasten 2: Konnektoren}

Beispiel: \es{Mi amiga es inteligente y \lsg{también} muy simpática, \lsg{pero} a veces está cansada \lsg{porque} trabaja mucho.}

\vspace{0.8em}

\textbf{Aufgabe 1: \es{ser/estar} mit Bedeutungsunterschied (6 P.)}
\begin{enumerate}[label=\alph*), leftmargin=2em, nosep]
    \item Mi profesor \lsg{es} muy aburrido. (= Er ist langweilig.)
    \item Hoy \lsg{estoy} aburrida porque no hay nada que hacer.
    \item Los estudiantes \lsg{están} listos para el examen. (= bereit)
    \item Tu hermana \lsg{es} muy lista. (= klug)
    \item Normalmente \lsg{soy} tranquilo, pero hoy \lsg{estoy} nervioso.
    \item ¿Dónde \lsg{está} María? -- \lsg{Está} en la biblioteca.
\end{enumerate}

\vspace{0.5em}

\textbf{Aufgabe 2: Adjektive und Antonyme (6 P.)}

\begin{tabularx}{\textwidth}{@{}|X|X|X|X|@{}}
\hline
simpático/-a & \lsg{sympathisch} & antipático/-a & \lsg{unsympathisch} \\
\hline
trabajador/-a & \lsg{fleißig} & perezoso/-a & \lsg{faul} \\
\hline
optimista & \lsg{optimistisch} & pesimista & \lsg{pessimistisch} \\
\hline
\end{tabularx}

\vspace{0.5em}

\textbf{Aufgabe 3: Konnektoren (6 P.)}
\begin{enumerate}[label=\alph*), leftmargin=2em, nosep]
    \item \lsg{Mi amigo es simpático y también (es) inteligente.}
    \item \lsg{Quiero ser médico porque me gusta ayudar a la gente.}
    \item \lsg{Es muy trabajadora, pero / sin embargo a veces está cansada.}
\end{enumerate}

\vspace{0.5em}

\textbf{Aufgabe 4: Berufe mit Begründung (8 P.)}
\begin{enumerate}[label=\alph*), leftmargin=2em, nosep]
    \item \lsg{Para ser profesor/-a hay que ser paciente y creativo/-a, ya que los estudiantes necesitan explicaciones claras.}
    \item \lsg{Para ser artista hay que ser creativo/-a y trabajador/-a, porque el arte requiere mucha práctica.}
    \item \lsg{Para ser bombero/-a hay que ser valiente y responsable, ya que es un trabajo peligroso.}
    \item \lsg{Para ser veterinario/-a hay que ser paciente y responsable, porque los animales no pueden hablar.}
\end{enumerate}

\vspace{0.5em}

\textbf{Aufgabe 5: Textproduktion (9 P.) -- Beispiellösung:}

\lsg{Mi compañera de clase se llama Ana. Tiene 18 años y es de Rostock. Es muy inteligente y también bastante trabajadora. Además, es simpática y siempre ayuda a los demás. Sin embargo, a veces está un poco nerviosa porque tiene muchos exámenes. Me gusta trabajar con ella porque es muy organizada y creativa.}

\vspace{0.3em}
\textit{\textcolor{grau}{Bewertung: 6 Sätze min. (3P), ser/estar korrekt (2P), 3 Konnektoren (2P), Steigerungen (2P)}}

\end{gruppeBbox}

\vspace{1em}
\hrule
\vspace{0.3em}

\textbf{Bewertungshinweise:}
\begin{itemize}[leftmargin=1.5em, nosep]
    \item Gruppe A: 35 Punkte gesamt
    \item Gruppe B: 35 Punkte gesamt
    \item Bei Textproduktion: Inhalt und Strukturvielfalt wichtiger als perfekte Grammatik
    \item Halbe Punkte möglich bei teilweise richtigen Antworten
    \item Aufgabe 4 (Interview): Punkte für verständliche, korrekte Sätze über Partner
\end{itemize}

\vspace{1em}

\textbf{Häufige Fehler und Tipps:}
\begin{itemize}[leftmargin=1.5em, nosep]
    \item \textbf{ser/estar-Verwechslung:} Immer fragen: Ist es eine Eigenschaft (ser) oder ein Zustand/Gefühl (estar)?
    \item \textbf{Adjektivendungen:} Auf Genus achten! \es{simpático} (m) vs. \es{simpática} (f)
    \item \textbf{hay que:} Immer mit Infinitiv, nie konjugiert!
    \item \textbf{Konnektoren:} \es{porque} (weil) vs. \es{pero} (aber) -- auf Bedeutung achten!
\end{itemize}

\end{document}
